\chapter{Operating, monitoring, and maintaining a server}
\label{ops:administration}

\section{A repeatable server configuration}

For a learning session, positional arguments are sufficient. For an ongoing
installation, copy \path{config/minisql.example.json}, change the paths and
policy deliberately, and start a configuration-driven listener. MiniSQL does
not search the working directory for a configuration file automatically: the
file is an explicit command argument.

\begin{lstlisting}
minisqld.exe --serve-config C:\MiniSQL-Data\db_ACTUAL_ID `
  C:\MiniSQL-Config\local.json
\end{lstlisting}

\begin{lstlisting}[language=bash]
minisqld --serve-config /srv/minisql/db_ACTUAL_ID \
  /etc/minisql/local.json
\end{lstlisting}

This trusted configuration mode must bind exactly \code{127.0.0.1}.
Use \code{--serve-tls-config} for the native TLS listener, with
\code{tls.enabled} set to true and a valid platform-specific certificate
reference. The \code{--serve-authenticated-config} mode uses password
authentication and MiniSQL's protected framing. The
\code{--serve-standby-config} mode is a trusted loopback standby listener.
Choosing one mode is part of the command; changing \code{tls.enabled}
does not transform \code{--serve-config} into the TLS listener.

Configuration is loaded and validated at startup. There is no documented
live-reload command. Restart with the revised file when changing server
settings, and keep the previous known-good file for rollback. A malformed
or unknown key is rejected instead of silently ignored.


Positional launchers do not automatically apply the example JSON's runtime
settings. They configure default logging but use the lower-level server
defaults: for example, its query-memory budget is 64 MiB and ordinary
positional concurrent listeners use a 5000 ms logical lock wait. Configured
startup passes the JSON query-memory value (128 MiB by default) and its
query-timeout value to lock-wait handling. Use the configured forms when the
precise operational limits matter, and inspect SHOW STATUS for active values.

\sourcecode{src/apps/minisqld/main.ml}{runConfiguredServer and mode-specific dispatch.}
\sourcecode{src/minisql/config/loader.ml}{Required keys, optional sections, and unknown-key rejection.}

\section{Complete configuration reference}
\index{configuration}

The following tables cover every key in the current model. Values are the
compiled defaults unless an example-file difference is explicitly noted.
Byte limits are decimal integers in JSON. A MiB is 1,048,576 bytes; a GiB is
1,073,741,824 bytes.

The required top-level sections are \code{configVersion}, \code{paths},
\code{server}, \code{runtime}, \code{databaseDefaults}, and \code{safety}.
The \code{logging}, \code{binlog}, and \code{tls} sections are optional;
omitting a whole section selects its compiled defaults. Once one of those
sections is present, its documented members are required. Recent resource
controls in the server/runtime sections have fallback defaults so older
version-1 configuration structures remain readable, provided their persisted
format settings are updated to the current required versions.

\subsection{Version and paths}

\begin{longtable}{p{.36\linewidth}p{.57\linewidth}}
\toprule
Key / default & Meaning and restrictions \\
\midrule
\endhead
\code{configVersion}: 1 & Configuration structure version; exactly 1. This
does not mean database format 1. \\
\code{paths.dataRoot}: \path{./data} & Nonempty path retained in the model.
The configured daemon opens the explicit database-path argument; it does not
derive or create that database from this key. \\
\code{paths.temporaryRoot}: \path{./tmp} & Nonempty configured path. Current
query execution places spill files under \path{<database>/tmp}; this setting
does not relocate that directory in the daemon execution path. \\
\code{paths.logDirectory}: \path{./logs} & Nonempty directory used for
ordinary file logs and SQL binlogs. It is active in configured startup. \\
\bottomrule
\end{longtable}

Relative paths are relative to the server's working directory, not
necessarily the JSON file's directory. In JSON, escape a Windows backslash,
for example:

\begin{lstlisting}
"logDirectory": "C:\\MiniSQL-Logs"
\end{lstlisting}

Forward slashes in a Windows path are often easier to read. Prefer absolute
deployment paths and a documented working directory.

\subsection{Server limits}
\index{connection limits}\index{frame limits}

All entries in the next table have prefix \code{server.}.

\begin{longtable}{p{.30\linewidth}p{.19\linewidth}p{.44\linewidth}}
\toprule
Key & Default & Effective operational range / meaning \\
\midrule
\endhead
\code{bindAddress} & \code{127.0.0.1} & Nonempty address. Trusted and standby
configured modes require the exact loopback address. Remote exposure needs
the intended authenticated transport and validation policy. \\
\code{port} & 7432 & Integer 1--65535. \\
\code{maxConnections} & 32 & Native concurrent listener accepts 1--128.
This bounds active workers/sessions, not write parallelism. \\
\code{maxStatementBytes} & 1048576 & At least 1024; maximum encoded SQL
request payload admitted by the database. Wire framing also has its own
protocol ceiling. \\
\code{maxFrameBytes} & 8388608 & 1024--16777216; encoded response payload
limit per frame. Large results use continuation frames. \\
\code{maxResultRows} & 1000000 & Positive integer; maximum returned rows per
statement across materialized and streamed delivery. \\
\code{maxResultBytes} & 268435456 & At least 1048576; cumulative encoded
response-byte limit per statement. \\
\code{idleTimeoutMs} & 300000 & At least 1000; inactive client lifetime in
milliseconds. \\
\bottomrule
\end{longtable}

\code{maxResultRows}, \code{maxResultBytes}, and \code{idleTimeoutMs} may
be omitted from an older server object and receive the defaults above. The
remaining server fields are required. The generic JSON validator initially
accepts some positive values more loosely, but the database/listener rejects
values outside the effective operational ranges. In particular, the schema's
unbounded positive \code{maxConnections} does not override the listener's
128-worker ceiling.

\subsection{Runtime budgets}
\index{memory budgets}\index{temporary storage}

All entries in the next table have prefix \code{runtime.}.

\begin{longtable}{p{.30\linewidth}p{.19\linewidth}p{.44\linewidth}}
\toprule
Key & Default & Meaning and validation \\
\midrule
\endhead
\code{bufferPoolBytes} & 268435456 & Shared read-cache budget. The JSON
schema recommends at least 1 MiB; the database open path requires at least
4096 and the actual page-cache construction must accommodate its page size. \\
\code{queryTimeoutMs} & 30000 & Positive absolute cooperative request
deadline, including waiting and streaming. Also supplied to logical lock-wait
handling by configured startup. \\
\code{checkpointWalBytes} & 67108864 & Threshold for WAL/checkpoint
maintenance. Schema minimum 1 MiB; database open path minimum 4096. \\
\code{temporaryMemoryBytes} & 134217728 & Soft per-session blocking-operator
budget; at least 1 MiB at database open. \\
\code{processMemoryBytes} & 2147483648 & Managed-heap admission/poll ceiling;
at least 16 MiB, and configuration validation requires it to cover the sum
of buffer-pool and temporary-memory settings. \\
\code{temporaryStorageBytes} & 1073741824 & Shared concurrent spill-reservation
limit; at least 1 MiB. \\
\code{slowQueryMs} & 1000 & Positive duration threshold for slow-query
counting and WARNING records. \\
\code{logLevel} & \code{info} & One of \code{debug}, \code{info},
\code{warning}, \code{warn}, or \code{error}. \code{warn} is an alias. \\
\bottomrule
\end{longtable}

The process-memory, temporary-storage, and slow-query settings are optional
members with these fallback values. Other runtime members are required.
Configuration validation and \path{config/config.schema.json} differ in a
few lower bounds. Use values satisfying the published schema and the effective
runtime checks; relying on a smaller value accepted by only one layer makes
configuration validation less portable.

\subsection{Logging and binlog}

\begin{longtable}{p{.36\linewidth}p{.57\linewidth}}
\toprule
Key / default & Meaning and restrictions \\
\midrule
\endhead
\code{logging.stdoutEnabled}: true & Boolean destination switch. \\
\code{logging.fileEnabled}: true & Boolean destination switch. At least one
ordinary destination must remain enabled. \\
\code{logging.fileName}: \path{minisql.log} & Nonempty file name inside the
log directory; no slash, backslash, colon, or NUL. \\
\code{logging.rotationHours}: 24 & Integer 1--87600; elapsed-time rotation
of a nonempty active log. \\
\code{binlog.enabled}: false & Enables complete valid UTF-8 SQL statement
capture independently of ordinary severity filtering. \\
\code{binlog.fileName}: \path{minisql-bin.log} & Nonempty file name with the
same filename restrictions. Must differ from the ordinary log filename. \\
\bottomrule
\end{longtable}

The singleton logger serializes records from multiple workers, preventing
their file writes from interleaving. It opens enabled destinations eagerly:
an invalid log directory fails startup instead of quietly losing later
records. File output is written before stdout. Rotation creates successive
archive files; it is not an automatic retention/deletion policy. Manage disk
capacity and retention outside the server.

\subsection{TLS configuration}

All entries in the next table have prefix \code{tls.}.

\begin{longtable}{p{.30\linewidth}p{.63\linewidth}}
\toprule
Key & Default and accepted meaning \\
\midrule
\endhead
\code{enabled} & false. Must be true for \code{--serve-tls-config}. \\
\code{certificateReference} & Compiled default \code{store:}; the example
file contains \path{pfx:.\\config\\server.pfx}. Neither is a ready-to-use
certificate. Set an actual \code{store:}, \code{pfx:}, or Linux
\code{pem:cert-path|key-path} reference. Nonempty when enabled. \\
\code{pfxPasswordEnvironment} & Exactly
\code{MINISQL\_TLS\_PFX\_PASSWORD}. \\
\code{protocolVersion} & Exactly \code{TLS1.3}. \\
\code{cipherSuite} & Exactly \code{TLS\_AES\_256\_GCM\_SHA384}. \\
\code{namedGroup} & Exactly \code{X25519}. \\
\bottomrule
\end{longtable}

The validation rule for a non-loopback bind requires either the TLS-enabled
flag or the explicit unsafe override in the safety model. That flag is a
configuration check; the selected command still determines the actual
listener. Use \code{--serve-tls-config} when intending TLS. Do not satisfy
validation by turning on a flag and then start a different transport mode.

\subsection{Persisted database defaults}

All entries below have prefix \code{databaseDefaults.}. They are required by
the JSON structure. Existing database format parameters are read from its
metadata and are not changed by startup defaults.

\begin{longtable}{p{.34\linewidth}p{.59\linewidth}}
\toprule
Key & Default and accepted values \\
\midrule
\endhead
\code{pageSize} & 4096; one of 4096, 8192, 16384, 32768. \\
\code{checksumAlgorithm} & Exactly \code{crc32c}. \\
\code{walSegmentBytes} & 16777216; runtime validator minimum 4096, JSON
schema minimum 1048576. Do not treat it as a complete archive-retention policy. \\
\code{textEncoding} & Exactly \code{utf-8}. \\
\code{defaultCollation} & Exactly \code{binary}. \\
\code{databaseFormatVersion} & Exactly 2. \\
\code{tableFileFormatVersion} & Exactly 2. \\
\code{indexFileFormatVersion} & Exactly 2. \\
\code{walFormatVersion} & Exactly 1. \\
\code{rowFormatVersion} & Exactly 2. \\
\bottomrule
\end{longtable}

\subsection{Safety settings}

\begin{longtable}{p{.42\linewidth}p{.51\linewidth}}
\toprule
Key & Meaning \\
\midrule
\endhead
\code{safety.durability} & Exactly \code{full}. There is no supported
relaxed durability mode. \\
\code{safety.allowUnknownFormatFeatures} & false; true is rejected.
Unknown persisted feature bits must not be accepted speculatively. \\
\code{safety.\allowbreak allowRemote\allowbreak WithoutAuthentication} & false in the example,
compiled default, and schema. The runtime model permits true as an explicit
unsafe override of the non-loopback configuration check. It does not change
the command's authentication mode, and trusted/standby configured listeners
still reject non-loopback addresses. Leave false in supported configurations. \\
\bottomrule
\end{longtable}

\sourcecode{config/minisql.example.json}{Complete example configuration.}
\sourcecode{config/config.schema.json}{Published schema and lower bounds.}
\sourcecode{src/minisql/config/model.ml}{Compiled defaults.}
\sourcecode{src/minisql/config/validation.ml}{Cross-field validation and safety policy.}
\sourcecode{src/minisql/server/database_manager.ml}{Effective open-time and operational limits.}
\sourcecode{src/minisql/server/listener.ml}{Concurrent listener worker ceiling.}

\section{Size memory, results, and temporary storage together}

The three memory-related settings serve different purposes. The buffer pool
is a shared cache. Temporary memory is a soft per-session operator budget.
Process memory is a managed-heap admission and cooperative-abort threshold.
The last is not an operating-system resident-set-size ceiling: native TLS
state, sockets, thread stacks, mappings, and other runtime allocations can
fall outside it. Apply the final process/container resource boundary at the
OS layer when one is required.

Hash joins, grouped aggregation, and sort use sampled row width to derive
thresholds and spill eligible work to the database's temporary directory.
A large row or an operator that still materializes its inputs can exceed a
soft query budget. The compatibility result API can also retain data beyond
an operator's internal budget. Streaming result frames limit transport chunks;
they do not prove that every query plan is constant-memory.

Increasing \code{maxConnections} while keeping a large per-query budget can
raise peak memory substantially. Increasing \code{maxResultBytes} can permit
very large result transfers even if each frame is small. Before changing a
limit, identify whether the observed failure is a row limit, response-byte
limit, spill limit, heap limit, or deadline. The correct remedy might be a
more selective query, projection of fewer columns, pagination, an index, or a
smaller batch rather than a larger global limit.

\section{Inspect status and active sessions}
\index{SHOW STATUS}\index{SHOW PROCESSLIST}

\begin{lstlisting}[language=SQL]
SHOW STATUS;
SHOW PROCESSLIST;
\end{lstlisting}

Authenticated callers require database ADMIN for both. Trusted sessions are
administrative. SHOW STATUS returns \code{variable\_name} and
\code{value} text columns; convert numeric values when graphing them. The
counters are process-local and reset on restart.

\begin{longtable}{p{.29\linewidth}p{.64\linewidth}}
\toprule
Group & Status names \\
\midrule
\endhead
Lifetime and work & \code{uptime\_ms}, \code{active\_sessions},
\code{total\_connections}, \code{total\_statements},
\code{failed\_statements}, \code{rows\_returned},
\code{checkpoint\_resets}. \\
Request/result policy & \code{max\_statement\_bytes},
\code{max\_frame\_bytes}, \code{max\_result\_rows},
\code{max\_result\_bytes}, \code{idle\_timeout\_ms},
\code{query\_timeout\_ms}. \\
Memory/spill & \code{query\_memory\_bytes}, \code{process\_memory\_bytes},
\code{heap\_used\_bytes}, \code{heap\_committed\_bytes},
\code{temporary\_storage\_bytes}, \code{temporary\_reserved\_bytes},
\code{temporary\_peak\_bytes}. \\
Failures and timing & \code{cancelled\_statements},
\code{timed\_out\_statements}, \code{resource\_rejected\_statements},
\code{total\_execution\_ms}, \code{maximum\_execution\_ms},
\code{slow\_query\_count}, \code{slow\_query\_ms},
\code{result\_bytes\_returned}. \\
Write fencing & \code{fencing\_enabled}, \code{fencing\_epoch},
\code{fencing\_rejections}. \\
\bottomrule
\end{longtable}

SHOW PROCESSLIST returns \code{session\_id}, \code{principal},
\code{peer}, \code{tls}, \code{state}, \code{state\_time\_ms},
\code{statement}, and \code{requests}. The statement field is a bounded
summary, not a complete audit trail. The TLS flag is set from the presence
of a native TLS credential on the listener; it does not describe the caller's
certificate trust policy. MiniSQL-specific protected framing can be in use
while this TLS flag is false.

Interpret age together with state. An idle connection is different from a
long-running query. A waiting query is different from one doing expensive
work. Keep a timestamped status snapshot when investigating a problem so a
later counter increase can be distinguished from activity before the incident.

\sourcecode{src/minisql/executor/executor.ml}{executeShowStatus and executeShowProcesslist.}
\sourcecode{src/minisql/server/listener.ml}{registerSessionPeer TLS argument.}

\section{Cancel work without abandoning all other clients}

Find the session ID with SHOW PROCESSLIST. Use the supported client API on a
separate administrative connection to request cancellation. Python exposes
\code{cancel\_session(id)} for an identified target and
\code{Connection.cancel()} for the connection's own active operation.
JDBC exposes \code{Statement.cancel()}. Native clients have
\code{cancelSession(client, id)}.

Cancellation is cooperative. It can take effect at an execution poll boundary,
and the client still needs to handle the resulting error and transaction
state. It is not a SQL \code{KILL} command or a guarantee that a thread has
been forcibly destroyed. If an explicit transaction is failed, roll it back
or recover through a valid savepoint before further ordinary work.

The control connection itself needs an available server connection slot and,
in authenticated mode, ADMIN. Leave operational capacity instead of consuming
all slots with an oversized application pool. A separate connection is used
because a blocked query socket cannot simultaneously carry an independent
administrative request.

\section{Monitor through the Prometheus bridge}
\index{Prometheus}\index{monitoring}

The repository includes a small exporter that uses the Python connector to
run SHOW STATUS and expose numeric values as Prometheus text:

\begin{lstlisting}[language=bash]
python3 tools/monitoring/minisql_exporter.py \
  --database-port 7432 --database learning --listen-port 9107
\end{lstlisting}

The default listen address is \code{127.0.0.1}; the exporter rejects
non-loopback binds. Both \path{/metrics} and \path{/healthz} perform a
database scrape. Success returns HTTP 200 with \code{minisql\_up 1};
failure returns HTTP 503 with \code{minisql\_up 0}. An exporter process
that is running but cannot reach its database must therefore be treated as an
unhealthy scrape, not healthy monitoring.

\begin{longtable}{p{.32\linewidth}p{.61\linewidth}}
\toprule
Exporter option & Default / purpose \\
\midrule
\endhead
\code{--listen-address} & \code{127.0.0.1}; loopback-only policy. \\
\code{--listen-port} & 9107. \\
\code{--database-host} & \code{127.0.0.1}. \\
\code{--database-port} & 7432. \\
\code{--database} & \code{main}; set the intended logical connection name. \\
\code{--user} & Omitted for trusted mode; use an ADMIN-capable authenticated
account otherwise. \\
\code{--password-env} & \code{MINISQL\_EXPORTER\_PASSWORD}; name of the
environment variable containing the secret. \\
\code{--tls} & Enables driver TLS. \\
\code{--server-name} & \code{localhost}; expected TLS identity. \\
\code{--pin-sha256} & Optional expected leaf-certificate digest. \\
\code{--ca-file} & Optional trusted CA file. \\
\code{--timeout} & 5.0 seconds for connect/socket operations. \\
\code{--self-test} & Runs the metrics-rendering check without serving HTTP. \\
\bottomrule
\end{longtable}

The current exporter renders all numeric status values with Prometheus
\code{gauge} type declarations, including cumulative server counters. Build
dashboards with their actual monotonic/reset behavior in mind. Authentication
and remote exposure of the HTTP endpoint belong in a protected reverse proxy;
the exporter does not provide its own remote HTTP authentication service.

\sourcecode{tools/monitoring/minisql_exporter.py}{Options, scrape behavior, metrics names, and HTTP health policy.}

\section{Routine maintenance: ANALYZE, REINDEX, and VACUUM}
\index{ANALYZE}\index{REINDEX}\index{VACUUM}

\begin{lstlisting}[language=SQL]
ANALYZE;
ANALYZE products;
REINDEX;
REINDEX products;
REINDEX product_name_index;
VACUUM;
VACUUM products;
\end{lstlisting}

ANALYZE refreshes optimizer statistics. Use it after substantial changes to
row counts or data distributions and inspect plans before attributing a slow
query to storage corruption. Statistics are advisory; they guide choices
without changing what the SQL means. The performance chapter explains their
sampled distributions and query-plan effects.

REINDEX rebuilds derived indexes for the whole database, a table, or an
explicit index. It does not reconstruct missing authoritative table rows.
VACUUM rewrites selected table storage, rebuilds indexes, and refreshes
statistics. It is useful after enough deletes/updates have changed physical
space use, but it has a greater I/O and temporary-space cost than simply
updating statistics.

All three require autocommit and database MAINTAIN for authenticated users.
They execute exclusively. Schedule the work with its expected duration and
free-space requirement in mind; do not start it inside an application DML
transaction. DDL/ANALYZE/VACUUM/REINDEX also invalidate relevant cached
planning generations so later statements can be planned against current
metadata.

VACUUM and operations that rewrite/reset physical state break the continuity
assumptions of a DML-only WAL archive. Stop the replication workflow and make
a new base after the maintenance or migration. A WAL archive is not an
independent copy of every schema and physical-layout transition.

\sourcecode{src/minisql/executor/executor.ml}{executeVacuum, executeReindex, and analyze handling.}
\sourcecode{src/minisql/executor/dml.ml}{Heap rewrites and index rebuild maintenance.}

\section{Offline consistency checking}
\index{consistency checker}

Stop the server, preserve a backup when investigating damage, and run:

\begin{lstlisting}
minisql-check.exe C:\MiniSQL-Data\db_ACTUAL_ID
\end{lstlisting}

A successful run prints a SUCCESS summary with table, row, and index counts,
and may print warnings. The checker streams heap rows, external TEXT/BLOB
chains, and active B+ tree leaves. It proves heap/index agreement using exact
entry probes and count agreement, so it is bounded by an individual row
rather than retaining the entire dataset. Its time still scales with the
authoritative data that must be read and verified.

The checker opens through the normal database manager, which can perform
recovery and derived repair as part of opening. Treat it as an offline
verification operation with an exclusive database lock, not a purely
read-only forensic acquisition utility. Preserve the original directory
before checking a copy when unchanged evidence is required.

A successful structural check does not prove that every application invariant
is satisfied. In particular, it is not a substitute for queries that look for
orphan relationships or for the documented SQL implementation caveats. Run
semantic checks relevant to the application after restore or migration.

Small \path{*.heap-pages} files beside table files are derived physical scan
directories. Missing, stale, or damaged copies rebuild from checksummed table
pages. Deleting one while the database is stopped can be a safe targeted
repair; copying one from another table or physical generation is not. Do not
extend this advice to table files, metadata, WAL, or index dirty markers.

\sourcecode{src/minisql/tools/check.ml}{Streaming validation and structural checks.}
\sourcecode{src/apps/minisql_check/main.ml}{Public checker command and exit behavior.}

\section{Service supervision and shutdown}
\index{SHUTDOWN}\index{systemd}\index{Task Scheduler}

\subsection{Linux systemd template}

The supplied \path{deploy/systemd/minisql.service} uses a dedicated
\code{minisql} user/group, a fixed working directory, a trusted configured
listener, restart-on-failure, and an ExecStop SQL SHUTDOWN command. Its default
paths are template values. Replace the database path with the actual
initialized directory, install the binaries/configuration where the unit
expects them, and create the service account and directories with appropriate
ownership.

The template's filesystem protection permits writes under its declared data
and log paths. If configuration selects a different write location, update
that policy as well. The stop command uses port 7432; changing the server's
port requires changing the stop command. For an authenticated/TLS listener,
replace the trusted stop command with an operator-managed credential helper
that issues SHUTDOWN through the matching transport.

\subsection{Windows templates}

\path{deploy/windows/Start-MiniSQL.ps1} accepts
\code{-DatabasePath}, \code{-ConfigPath}, and optional
\code{-BinaryDirectory}. It resolves paths and runs the trusted configured
daemon in the foreground. \path{Stop-MiniSQL.ps1} accepts \code{-Port}
and the binary directory and sends SHUTDOWN through the trusted client.
Use them with Task Scheduler or a suitable supervisor if unattended operation
is needed.

These scripts do not register a native Windows Service Control Manager
service. The current MiniLang runtime does not expose the SCM callback ABI
needed to claim that integration. A wrapper can supervise the process, but
the wrapper's startup, account, working directory, logging, and graceful-stop
behavior must be configured explicitly.

\subsection{Graceful and abnormal stops}

SQL SHUTDOWN acknowledges the request and asks the listener to drain workers
and close storage cooperatively. The console also advertises Ctrl+C for
stopping an operational server. Prefer the SQL path for planned maintenance
when an administrative connection is available. After an abnormal process
termination, open/recovery and an offline check provide evidence before the
database is returned to important work.

Do not start a second writable server against the same database directory.
The database lock exists to prevent conflicting ownership. Deleting its file
while a process still owns storage does not create safe multiwriter access.

\sourcecode{deploy/systemd/minisql.service}{Linux supervision example.}
\sourcecode{deploy/windows/Start-MiniSQL.ps1}{Windows startup wrapper.}
\sourcecode{deploy/windows/Stop-MiniSQL.ps1}{Trusted shutdown wrapper.}

\section{Daemon command reference}

The forms below omit \path{.exe} for readability. Angle-bracket names are
placeholders, not literal path text. Optional \code{max-clients} defaults to
32 for the ordinary persistent modes; concurrent modes accept at most 128.
An optional request budget of zero means unbounded operation.

\begin{lstlisting}
minisqld --version
minisqld --m0-self-test
minisqld --init <data-root> <name> [page-size]
minisqld --set-admin-password <database>
minisqld --set-user-password <database> <user>
minisqld --generate-encryption-key <new-key-file>
minisqld --enable-encryption <database> <key-file>
minisqld --rotate-encryption-key <database> <new-key-file>

minisqld --serve <database> <port> [max-clients]
minisqld --serve-authenticated <database> <port> [max-clients]
minisqld --serve-standby <standby> <port> [max-clients]
minisqld --serve-secure <database> <address> <port>
  <max-clients> [max-requests]
minisqld --serve-tls <database> <address> <port>
  <max-clients> <certificate-reference> [max-requests]
minisqld --serve-fenced <database> <port> <max-clients>
  <lease-file> <epoch> <node-id> <max-clock-skew-ms>

minisqld --serve-config <database> <config-file>
minisqld --serve-authenticated-config <database> <config-file>
minisqld --serve-standby-config <standby> <config-file>
minisqld --serve-tls-config <database> <config-file>

minisqld --serve-one <database> <port> [max-requests]
minisqld --serve-auth <database> <port> [max-requests]
minisqld --serve-many <database> <port> <max-clients>
  [max-requests]
minisqld --serve-auth-many <database> <port> <max-clients>
  [max-requests]
\end{lstlisting}

Wrapped lines in this syntax reference are one command. The one-session and
request-budget forms are retained for compatibility and tests. Use persistent
operational forms for a service. The \code{--m0-self-test} banner is a
historical interface check, not the full test suite. Usage errors generally
return exit code 2; an operational error prints \code{ERROR number: message}
and returns 1; successful commands return 0. Inspect both the exit status and
output in administrative automation.

\sourcecode{src/apps/minisqld/main.ml}{Complete public daemon dispatch.}

\section{Operational qualification and incident records}

The repository includes workload, capacity, fault, recovery, and HA drills.
Run destructive drills only in their intended new disposable work roots.
The production fault drill tests controlled storage exhaustion, native crash
boundaries, interrupted TCP frames, termination during concurrent commits,
restart, continued writes, and corruption rejection. It deliberately avoids
filling the host volume. A real filesystem-full test belongs on a disposable
quota-limited volume.

The soak script measures a representative persistent workload over an extended
period and can fail on workload errors, process disappearance, RSS/handle/
thread drift, or latency thresholds. Historical results are evidence about
the tested revision, hardware, and workload. Preserve a report for the
deployment's actual binary and largest representative database.

When an incident occurs, record the source revision, binary versions,
configuration, SQL/error text with secrets removed, time, process status,
disk free space, status counters, and relevant logs. Before attempting repair,
copy the complete stopped directory and its required keys to a controlled
location. Prefer restoring a verified copy over modifying the only surviving
data in place.

\sourcecode{tests/fault/production_fault_drill.py}{Isolated fault qualification.}
\sourcecode{tests/soak/production_soak.py}{Long-running workload/resource checks.}
\sourcecode{tests/integration/production_controls_live.py}{Operational limit and cancellation scenarios.}
